\bookStart{Galder of Growe}[Gróu galdr]

% Introduction.
\begin{flushright}%
\textbf{Dating} \parencite{Robinson1991}[397--406]: 1200--1250.

\textbf{Meter:} \Fornyrdislag%
\end{flushright}

\section{Introduction}

The \textbf{Galder of Growe} is a short mythological TODO: Introduction.

\subsection{Textual preservation and transmission}

\Grougaldr\ is preserved together with \textlink{Fjolsvinnsmal} in at least 46 post-reformation paper manuscripts.  As shown by \textcite{Robinson1991}, the two main manuscripts are Rask 21 a (\textbf{Ra}) and Stockholm Papp. 15 8vo (\textbf{St}).

\subsection{Dating}

Since \Grougaldr\ and \textlink{Fjolsvinnsmal} are related compositions, it is best to discuss the dating of both poems in the same place.  To begin with, we may try to obtain a \emph{terminus ante quem} for the two texts.

Both \Grougaldr\ and \textlink{Fjolsvinnsmal} are composed in a classical \Ljodahattr\ meter, which in both poems consistently observes Bugge’s rule, namely that a \Ljodahattr\ c-verse (or full line) cannot end in a dip.  In practice, this means that the c-verse can never end in a metrical trochée (/ x), e.g., the disyllabic word \emph{víti}.  On other hand, disyllabic words with light stems like \emph{viti} or \emph{saman} are allowed, since they resolve into a single lift.  This distinction could only have been relevant before the Icelandic lengthening of vowels in open syllables (i.e., CVCV > CVVCV), which would have made words like \emph{víti} and \emph{viti} metrically identical; a poet working after this sound change can not be expected to have been aware of the historical distinction in syllable quantity.  Since this lengthening was complete by ca.~1400 \ce, the poems must at the very least be mediaeval \parencite[401]{Robinson1991}.

\sectionline

\section{Text}

...TODO...

\bvg\bva%
Þann gęl’k þér inn \alst{f}immta, \hld\ ef þér \alst{f}jǫturr verðr &
\ind \alst{b}orinn at \alst{b}óg-limum: &
\alst{L}ęifnis ęlda \hld\ lę́t’k þér fyr \alst{l}ęgg of kveðinn, &
\ind ok støkkr þá \alst{l}ȧss af \alst{l}imum, &
\ind ęn af \alst{f}ótum \alst{f}jǫturr.\eva

\bvb This one I gale thee the fifth, if a fetter be \\
\ind borne onto thy shoulder-limbs: \\
Lafner’s fires I make to be sung about thy leg \\
\ind and then springs the lock from the limbs, \\
\ind and from the feet the fetter.\evb\evg


\bvg\bva%
Þann gęl’k þér inn \alst{s}étta, \hld\ ef þú ȧ \alst{s}jó kømr &
\ind \alst{m}ęira an \alst{m}ęnn viti: &
\alst{l}ogn ok \alst{l}ǫgr \hld\ gangi þér i̇́ \alst{l}úðr saman &
\ind ok ljái þér ę́ \alst{f}rið-drjúgrar \alst{f}arar.\eva

\bvb This one I gale thee the sixth, if thou comest to a sea \\
\ind greater than men might know:
may the calm and the water go together into the mill for thee, \\
\ind and always give thee a lasting peaceful journey.\evb\evg

...TODO...

\sectionline
